\section{Podsumowanie}
\label{sec:podsumowanie}

\subsection{Wnioski}

\begin{enumerate}
    \item Postawiony problem łączy aspekty kombinatoryczne (wybór podzbioru
          atrakcji z~ograniczeniami typów i~budżetów) z~grafowymi
          (istnienie i~koszt fizycznej ścieżki na siatce 8-kierunkowej
          z~zakazem powtórzeń krawędzi). Tradycyjne sformułowania
          \emph{Orienteering Problem} stanowią dobrą inspirację, ale
          twarde wymaganie braku powtórzeń krawędzi czyni nasz model
          wyraźnie trudniejszym.
    \item Dwupoziomowa reprezentacja rozwiązania (waypointy + materializacja
          A\textsuperscript{*}) okazała się praktyczna --- pozwala
          rozdzielić problem ,,co odwiedzić'' od~,,jak dojść''
          i~drastycznie zawęża przestrzeń sąsiedztwa.
    \item Algorytm ABC daje populację rozwiązań w~pełni dopuszczalnych,
          co eliminuje konieczność funkcji kary i~upraszcza analizę.
    \item Leksykograficzna krotka porządkowa
          ($\mathrm{Value}\uparrow$, $\mathrm{Time}\downarrow$,
          $\mathrm{Cost}_{\mathrm{ruch}}\downarrow$,
          $\mathrm{Cost}_{\mathrm{atrakcji}}\downarrow$) jest tanim, ale
          skutecznym kryterium różnicowania rozwiązań i~kierowania
          eksploatacją.
\end{enumerate}

\subsection{Napotkane problemy}

\begin{itemize}
    \item \textbf{Trudność generacji rozwiązań dopuszczalnych.} W~instancjach
          o~ciasnym budżecie czasowym lub bardzo dużej liczbie atrakcji
          częstość uzyskiwania ścieżek dopuszczalnych przy losowych
          waypointach jest mała --- konieczne jest podbicie liczby prób
          (\texttt{max\_tries}) w~\texttt{\_random\_feasible\_solution}.
    \item \textbf{Twardy zakaz powtórzeń krawędzi.} Wprowadzenie tego
          ograniczenia powoduje, że pewne kombinacje waypointów po prostu
          nie da się fizycznie zrealizować (segment domyka się dopiero
          po długim ,,objeździe''); musi to być uwzględnione w~mutacjach
          jako naturalny kanał odrzuceń.
    \item \textbf{Koszt A\textsuperscript{*} przy każdej mutacji.}
          Każda zaakceptowana mutacja oznacza materializację co najmniej
          jednego segmentu ścieżki; przy populacji 50 i~200 iteracjach
          koszt ten kumuluje się szybko --- optymalizacja A\textsuperscript{*}
          (np.\ z~cachowaniem segmentów lub heurystyką niższej jakości,
          ale szybszą) jest naturalnym kierunkiem przyspieszania.
\end{itemize}

\subsection{Kierunki rozwoju}

\begin{enumerate}
    \item \textbf{Pełna integracja ABC z~aplikacją Streamlit.}
          Architektura \texttt{app.py} przewiduje stabilny interfejs
          \texttt{solve(map, params) → path}; dorzucenie wrappera nad
          \texttt{generate\_abc\_population} pozwoli porównywać ABC
          z~baselineami w~UI.
    \item \textbf{Drugi algorytm meta-heurystyczny.} GA / SA / Tabu Search
          na tej samej reprezentacji waypointów (operatory już mamy)
          dałyby porównanie wewnątrz tej samej klasy metod.
    \item \textbf{Solver ILP referencyjny.} Dla małych instancji
          (np.\ $10 \times 10$, 6--8 atrakcji) sformułowanie z~rozdziału
          \ref{sec:model} można zlecić solverowi --- daje ,,prawdziwie
          optymalną'' wartość celu i~pozwala obiektywnie ocenić jakość
          ABC.
    \item \textbf{Bardziej wyrafinowana funkcja oceny.} Wprowadzenie
          metryki Pareto-typu (niezdominowana wartość vs.\ czas) zamiast
          leksykografii, by jawnie zwracać front Pareto.
    \item \textbf{Cache segmentów A\textsuperscript{*}.} Ponieważ ABC
          modyfikuje pojedyncze pozycje waypointów, większość segmentów
          ścieżki pozostaje niezmienna --- ich przeliczanie można
          zaoszczędzić.
    \item \textbf{Większa baza scenariuszy testowych.} Aktualnie dostępne
          są dwa scenariusze webowe (\texttt{small.json},
          \texttt{medium.json}); benchmark wymaga większej i~bardziej
          zróżnicowanej kolekcji.
\end{enumerate}
